\emph{Conventions.} Complex scalars are $\mathrm{span}\{1, I_{7}\}$, $I_{7} = I_{6}e$ central in $\Cl(3,4) = \Cl(2,4)_{\C}$ (\claimref{std-algebra.complexification-is-the-odd-extension}). $\CS$, $h$, $J = I_{7}e_{1}e_{+}$ and $\Sigma = h(\cdot, J\cdot)$ are those of Theorem \claimref{synthesis.positive-structure-is-twistor-form-at-imaginary-time}; in its model the generators are realised on three fermion modes in the Jordan-Wigner basis and $h$ is the identity Gram matrix. The check uses that model in exact Gaussian rationals $\mathrm{span}_{\mathbb{Q}}\{1, I_{7}\}$. Adjoints are antilinear in complex scalars.

\emph{Step 1 (the $\Sigma$-adjoint; proved).} By Theorem \claimref{synthesis.positive-structure-is-twistor-form-at-imaginary-time}(i), $J^{2} = 1$ and $J^{\dagger} = J$. For operators $A$, $\Sigma(\psi, A\varphi) = h(\psi, JAJ\,J\varphi) = h((JAJ)^{\dagger}\psi, J\varphi) = \Sigma(JA^{\dagger}J\psi, \varphi)$, so $A^{\sharp} = JA^{\dagger}J$. Replacing $\Sigma$ by $\lambda\Sigma$ with $\lambda$ real and nonzero multiplies both sides of the defining identity $\Sigma(A^{\sharp}\psi, \varphi) = \Sigma(\psi, A\varphi)$ by $\lambda$, so $A^{\sharp}$ is unchanged (the check also computes $(\lambda\Sigma)^{-1}A^{\dagger}(\lambda\Sigma)$ for $\lambda = -3, \tfrac27$). The map $A \mapsto A^{\sharp}$ is an antilinear anti-automorphism of $\operatorname{End}\CS$. By Theorem \claimref{synthesis.positive-structure-is-twistor-form-at-imaginary-time}(iii) the $\Sigma$-adjoint of every real vector is minus itself (directly: $e_{k}^{\dagger} = e_{k}^{-1} = e_{k}/Q(e_{k})$, and $J$ anticommutes with $e_{1}, e_{+}$, the generators of square $+1$, and commutes with the others, so $e_{k}^{\sharp} = -e_{k}$ for all six). Clifford conjugation is the unique anti-automorphism of $\Cl(2,4)$ that negates vectors, since the vectors generate the algebra; hence $A^{\sharp} = \bar A$ on $\Cl(2,4)$, and on a blade of grade $k$, $\bar B = (-1)^{k}\widetilde B = (-1)^{k}(-1)^{k(k-1)/2}B = (-1)^{k(k+1)/2}B$. The check compares $JB^{\dagger}J$ with $(-1)^{k(k+1)/2}B$ on all $64$ blades. This is (i).

\emph{Step 2 (trace forms on vectors; proved).} The complexified $\Cl(2,4)$ acts irreducibly on the eight-dimensional $\CS$, and $\tau(B) = 0$ for every blade $B \neq 1$: if $B$ has even grade below six it contains a generator $e$ with $eBe^{-1} = -B$; if $B$ has odd grade some generator $f$ is missing from $B$ and $fBf^{-1} = -B$; the pseudoscalar anticommutes with every generator. So $\tau(A) = \langle A\rangle_{0}$ for $A$ in the algebra. Then $\tau(a^{\sharp}b) = \langle\bar ab\rangle_{0} = -\langle ab\rangle_{0} = -a\cdot b$ for real vectors, a form of signature $(4,2)$ because $a\cdot b$ has signature $(2,4)$; and $\tau(e_{j}^{\dagger}e_{k}) = \langle e_{j}^{-1}e_{k}\rangle_{0} = \delta_{jk}$, the Euclidean form of the frame. By Step 1, $\tau(a^{\sharp}b) = \tau(Ja^{\dagger}Jb) = \tau(a^{\dagger}JbJ)$ by cyclicity, and $JbJ = (e_{1}e_{+})b(e_{1}e_{+})^{-1}$ ($I_{7}^{2} = -1$ and $(e_{1}e_{+})^{2} = -1$), which reverses $e_{1}$ and $e_{+}$ and fixes the other four generators. Restricted to $u, w$ in the span of $e_{1}, \dots, e_{4}$: $\tau(u^{\sharp}w) = -u\cdot w$, negative on $e_{1}$ and positive on $e_{2}, e_{3}, e_{4}$ (Gram matrix $\operatorname{diag}(-1, 1, 1, 1)$, checked), $\tau(u^{\dagger}w)$ is the Euclidean form with $e_{1}$ as time, and $JuJ = -e_{1}ue_{1}$ there. The check computes all $36$ Gram entries of both forms, the identity on three sample vectors with all six components, and the action of $J$. The identity $\langle\bar ab\rangle_{0} = -\langle ab\rangle_{0} = -a\cdot b$ uses only $\bar a = -a$, so it holds in every Clifford algebra; the check also computes it directly in $\Cl(1,3)$ and $\Cl(2,4)$ with no module and no twistor form. The sign that makes $e_{2}, e_{3}, e_{4}$ positive is the store's signature, $Q(e_{1}) = 1$; the normalisation is the choice of $\tau$ and of the sign in front of it. What the twistor form contributes is the identification of $\sharp$ with Clifford conjugation and the involution $J$ between the two forms. This is (ii).

\emph{Step 3 (invariance; proved).} For a rotor $R$ (even, $R\widetilde R = 1$), Clifford conjugation equals reversion on even elements, so $\overline{RaR^{-1}} = \overline{R^{-1}}\,\bar a\,\bar R = R\bar a\widetilde R$, and $\tau(\overline{RaR^{-1}}\,RbR^{-1}) = \langle R\bar ab\widetilde R\rangle_{0} = \langle\bar ab\rangle_{0}$. The check computes the invariance for the boost $\tfrac54 + \tfrac34e_{1}e_{2}$, the rotation $\tfrac35 + \tfrac45e_{2}e_{3}$, the translation $1 + \tfrac12e_{3}n_{\infty}$ ($n_{\infty} = e_{+} - e_{-}$) and the dilation $\tfrac54 + \tfrac34e_{+}e_{-}$, on sample vectors, by conjugating the $8 \times 8$ matrices. The boost takes $e_{1}$ to $\tfrac{17}{8}e_{1} \pm \tfrac{15}{8}e_{2}$, whose positive norm is $\tfrac{289 + 225}{64} = \tfrac{257}{32} \neq 1$ while its $\Sigma$-norm stays $-1$. That no positive structure is boost invariant is Theorem \claimref{synthesis.positive-structure-is-twistor-form-at-imaginary-time}(iv). This is (iii).

\emph{Step 4 (uniqueness; proved).} Let $X_{ij}$ be the matrices of $u \mapsto e_{i}e_{j}u - ue_{i}e_{j}$ on the span of $e_{1}, \dots, e_{4}$, computed in $\Cl(1,3)$. A real bilinear form with matrix $M$ is invariant under the Lorentz algebra if and only if $X_{ij}^{T}M + MX_{ij} = 0$ for the six pairs. These are $96$ linear equations in the $16$ entries of $M$; their solution space, computed exactly, is one-dimensional and spanned by $\operatorname{diag}(1, -1, -1, -1)$. A Hermitian form $H = S + I_{7}A$ ($S$ real symmetric, $A$ real antisymmetric) on $\mathrm{span}\{1, I_{7}\}^{4}$ is invariant under the real matrices $X_{ij}$ exactly when $S$ and $A$ are (the equations are real-linear and the $X_{ij}$ are real); $A$ lies in the solution space and is antisymmetric, so $A = 0$ and $H$ is a real multiple of the metric, hence of $-u\cdot w$, and not definite. This is the computed form of Schur's lemma for the real irreducible vector representation. This is (iv).

\emph{Step 5 (second quantisation; proved).} For a one-particle space with an orthogonal basis $z_{\mu}$ of norms $\kappa_{\mu} = \pm1$, the bosonic Fock form on monomials is $\langle z^{\alpha}, z^{\beta}\rangle_{K} = \delta_{\alpha\beta}\,\alpha!\prod_{\mu}\kappa_{\mu}^{\alpha_{\mu}}$. By Step 2, $\kappa = (-1, 1, 1, 1)$ for $\tau(u^{\sharp}w)$, so $\langle z^{\alpha}, z^{\beta}\rangle_{K} = (-1)^{\alpha_{0}}\alpha!\,\delta_{\alpha\beta} = \langle z^{\alpha}, (-1)^{n_{0}}z^{\beta}\rangle$ with the form $\langle z^{n}, z^{n'}\rangle = n!\,\delta$ of Proposition \claimref{field-theory.fock-space-maxwell-modes}. For $a_{\mu} = \partial/\partial z_{\mu}$: $\langle z^{\alpha}, a_{\mu}z^{\beta}\rangle_{K} = \beta_{\mu}\langle z^{\alpha}, z^{\beta - e_{\mu}}\rangle_{K}$, nonzero only for $\alpha = \beta - e_{\mu}$, where it is $\beta!\,\kappa^{\beta}/\kappa_{\mu} = \kappa_{\mu}\langle z_{\mu}z^{\alpha}, z^{\beta}\rangle_{K}$. So the $K$-adjoint of $a_{\mu}$ is $\kappa_{\mu}z_{\mu} = -g_{\mu\mu}z_{\mu}$ and $[a(u), a(w)^{\dagger}] = -u\cdot w$ for real polarisations, the form (K1) of Premise \claimref{synthesis.krein-form-premise-indefinite-inner-product-for-gauge-fields}; its fundamental symmetry $(-1)^{n_{0}}$ is the second quantisation of $u \mapsto JuJ$, which reverses $e_{1}$ only. The check verifies both identities on all monomials of degree at most three. The commutator $-u\cdot w$ on all four polarisations of each mode, with propagator $-g_{\mu\nu}/k^{2}$, is the Feynman-gauge ($\alpha = 1$) realisation. A covariant gauge with $\alpha \neq 1$ adds to the propagator a term $(1 - \alpha)k_{\mu}k_{\nu}/k^{4}$, whose two-point function has dipole support ($\delta'(k^{2})$); its state space needs the auxiliary (Nakanishi-Lautrup) field and is not a Fock space over a one-particle form at fixed momentum, so it is neither obtained nor excluded here (standard, cited from memory). This is (v).

\emph{Step 6 (the ghost twist; computed).} With $c = \tfrac12(e_{1} + I_{7}e_{+})$: $c^{2} = \tfrac14(e_{1}^{2} - e_{+}^{2}) = 0$ and $c^{\dagger}c = \tfrac14(e_{1} - I_{7}e_{+})(e_{1} + I_{7}e_{+}) = \tfrac12(1 + I_{7}e_{1}e_{+}) = \tfrac12(1 + J)$, so $n_{c} = c^{\dagger}c$ is a projection and $J = 2n_{c} - 1$: up to sign, $J$ is the parity of the mode $c$. On a floor of two tower modes $b, d$ with the CAR of Fact \claimref{field-theory.imported-tower-limit}, $n_{-} = \tfrac12(b^{*} - d^{*})(b - d)$ is the number operator of $(b - d)/\sqrt2$, $n_{-}^{2} = n_{-}$, and $S = 1 - 2n_{-}$, with $S^{2} = 1$, $S^{*} = S$ and $SbS = d$. Every fundamental symmetry that is the parity of one mode has this form, so the computation records the shape of (K2) and does not select it; the operative content of (K2) is the pairing of ghost and antighost in the indefinite form. This is (vi).

\emph{Step 7 (what does not follow; argued).} (a) The forms above are the covariant ones; the positive ones $\tau(a^{\dagger}b)$ and $h$ are equally present in the attachment. That a field theory's polarisation and ghost state spaces carry the covariant form rather than a frame's positive form is the content of (K1) and (K2) as a premise, and it is not derived here; Step 4 shows only that a Lorentz-invariant polarisation form is $-u\cdot w$ up to a real factor, which is standard. (b) Clause (K3) concerns additional massive fields entering loops with negative weights; nothing here says which fields exist. (c) Clause (K4) needs a BRST charge, that is, dynamics; positivity of its cohomology is computed in Theorem \claimref{synthesis.self-adjoint-brst-charge-needs-a-krein-form-and-the-tower-supplies-the-ghosts} only for the free charge at a null momentum up to three quanta. (d) The Fock form over an indefinite one-particle form is the standard construction; Clifford conjugation and its trace form, the Gupta-Bleuler form and the twistor form are standard. What is recorded is that, inside the store, the one-particle form of the premise is the trace form of the twistor adjoint (as of Clifford conjugation in any Clifford algebra), and that the fundamental symmetry and the frame of the positive form are supplied by the same $J$.

\emph{Status.} (i) to (v) proved, with exact checks; (vi) computed; the reading argued. A restatement of the one-particle content of (K1), not a derivation of the premise. Transport-absent.
